% Section content template for individual Educative lessons
% This file contains the actual content for a specific section
% Generated from Educative API section content

% Section: Code for the Hotel Management System
% Section ID: 4898729440313344
% Section Slug: code-for-the-hotel-management-system
% Generated: 2025-12-12T16:23:46.017668

We've reviewed different aspects of the hotel management system and observed the attributes attached to the problem using various UML diagrams. Let 's explore the more practical side of things, where we will work on implementing the hotel management system using multiple languages. This is usually the last step in an object-oriented design interview process.
\\
We have chosen the following languages to write the skeleton code of the different classes present in the hotel management system:

\begin{itemize}
\item
\protect\phantomsection\label{lfBIAqHx_QnUUW7JZlcQ0}
Java
\item
\protect\phantomsection\label{RCVAXLIMq9kTw5dUhVHNq}
C\#
\item
\protect\phantomsection\label{mDHLq2rjFXGCss_2aCDT4}
Python
\item
\protect\phantomsection\label{eU8ikc6Rp3zwNvtToX4O7}
C++
\item
\protect\phantomsection\label{lSqYOUUS7-9j4Kt8CfXV_}
JavaScript
\end{itemize}
\\
If you want to skip directly to the implementation and see the complete workflow, simply scroll down or click~\hyperref[Executable-code-Hotel-management-system]{Executable Code: Hotel Management System}.

\subsection{Hotel management system classes}\label{guqoYLhLV_iZhJI9I5gTP}
\\
This section will provide the skeleton code of the classes designed in the class diagram lesson.
\begin{quote}
\textbf{Note:} For simplicity, we are not defining getter and setter functions. The reader can assume that all class attributes are private, accessed through their respective public getter methods, and modified only through their public methods.
\end{quote}
\subsubsection{Enumerations}\label{mMjnqIP0mInF9wBuqDk2B}
\\
First, we will define all the enumerations required in the hotel management system. According to the class diagram, six enumerations are used in the system, i.e., \texttt{RoomStyle}, \texttt{RoomStatus}, \texttt{BookingStatus}, \texttt{AccountStatus}, \texttt{AccountType}, and \texttt{PaymentStatus}. The code to implement these enumerations is as follows:
\begin{quote}
\textbf{Note:} JavaScript does not support enumerations, so we will use the\ Object.freeze()\ method as an alternative that freezes an object and prevents further modifications.
\end{quote}

\textbf{Enum definitions}

\begin{lstlisting}[language=Java]
// definition of enumerations used in hotel management system
enum RoomStyle {
    STANDARD, 
    DELUXE, 
    FAMILY_SUITE, 
    BUSINESS_SUITE
}

enum RoomStatus {
    AVAILABLE, 
    RESERVED, 
    OCCUPIED, 
    NOT_AVAILABLE, 
    BEING_SERVICED, 
    OTHER
}

enum BookingStatus {
    REQUESTED, 
    PENDING, 
    CONFIRMED,
    CANCELLED, 
    ABANDONED
}

enum AccountStatus {
    ACTIVE, 
    CLOSED, 
    CANCELED, 
    BLACKLISTED, 
    BLOCKED
}

enum AccountType {
    MEMBER, 
    GUEST, 
    MANAGER, 
    RECEPTIONIST
}

enum PaymentStatus {
    UNPAID, 
    PENDING, 
    COMPLETED, 
    FILLED, 
    DECLINED, 
    CANCELLED, 
    ABANDONED, 
    SETTLING, 
    SETTLED, 
    REFUNDED
}
\end{lstlisting}

\subsubsection{Address and account}\label{RWXA6Mmrdlki4DsobSWvj}
\\
This section contains the \texttt{Address} and \texttt{Account} classes. Here , the \texttt{Address} class is used as a custom data type. The implementation of these classes is shown below:

\textbf{The Address and Account classes}

\begin{lstlisting}[language=Java]
public class Address {
    private String streetAddress;
    private String city;
    private String state;
    private int zipCode;
    private String country;
}

public class Account {
    private String id;
    private String password;
    private AccountStatus status;
  
    public boolean resetPassword();
}
\end{lstlisting}

\subsubsection{Person}\label{tDLcPrXhBdyCmy-YQ43tL}
\\
\texttt{Person} is an abstract class that represents the various people or actors that can interact with the system. There are four types of persons: \texttt{Guest}, \texttt{Receptionist}, \texttt{Server}, and \texttt{Housekeeper}. The implementation of the mentioned classes is shown below:

\textbf{Person and its child classes}

\begin{lstlisting}[language=Java]
public abstract class Person {
    private String name;
    private Address address;
    private String email;
    private String phone;
    private Account account;
}


public class Guest extends Person {
    private int totalRoomsCheckedIn;

    public List<RoomBooking> getBookings();
    public boolean createBooking();
}

public class Receptionist extends Person {
    public List<Person> searchMember(String name);
    public boolean createBooking();
}

public class Housekeeper extends Person {
    public boolean assignToRoom();
}
\end{lstlisting}

\subsubsection{Service}\label{LAMVTngPU8FxxoMdX-5bw}
\\
\texttt{Service} is an abstract class, and this section represents different charges (\texttt{Amenity}, \texttt{RoomService}, and \texttt{KitchenService}) against a booking. The code to implement these classes is shown below:

\textbf{Service and its derived classes}

\begin{lstlisting}[language=Java]
public abstract class Service {
    private Date issueAt;

    public boolean addInvoiceItem(Invoice invoice);
}

public class Amenity extends Service {
    private String name;
    private String description;
}

public class RoomService extends Service {
    private boolean isChargeable;
    private Date requestTime;
}

public class KitchenService extends Service {
    private String description;
}
\end{lstlisting}

\subsubsection{Invoice}\label{RM2PViZQ0YfzhcvZAiR0K}
\\
This section contains information on the \texttt{Invoice} class to create a bill. The implementation of this class is represented below:

\textbf{The Invoice class}

\begin{lstlisting}[language=Java]
public class Invoice {
    private double amount;
    
    public boolean createBill();
}
\end{lstlisting}

\subsubsection{Room booking}\label{SRSE9ivjxp87WdBSivLc7}
\\
\texttt{RoomBooking} is a class responsible for managing the bookings for a room. The implementation of this class is given below:

\textbf{The RoomBooking class}

\begin{lstlisting}[language=Java]
public class RoomBooking {
    private String reservationNumber;
    private Date startDate;
    private int durationInDays;
    private BookingStatus status;
    private Date checkin;
    private Date checkout;

    private int guestId;
    private Room room;
    private Invoice invoice;
    private List<Notification> notifications;

    public static RoomBooking fectchDetails(String reservationNumber);
}
\end{lstlisting}

\subsubsection{Bill transaction}\label{DpNAhdkSnHg7DcMraAzeH-1}
\\
After generating an invoice, a customer must pay the bill to confirm the room booking. A \texttt{BillTransaction} class is required to store the information about bill payment. Three ways to pay the bill are check, cash, and credit card transactions.

\textbf{BillTransaction and its derived classes}

\begin{lstlisting}[language=Java]
// BillTransaction is an abstract class
public abstract class BillTransaction {
    private Date creationDate;
    private double amount;      
    private PaymentStatus status;

    public abstract void initiateTransaction();
}

class CheckTransaction extends BillTransaction {
    private String bankName;
    private String checkNumber;

    public void initiateTransaction() {
        // functionality 
    }
}

class CreditCardTransaction extends BillTransaction {
    private String nameOnCard;
    private int zipcode;

    public void initiateTransaction() {
        // functionality 
    }
}

class CashTransaction extends BillTransaction {
    private double cashTendered;

    public void initiateTransaction() {
        // functionality 
    }
}
\end{lstlisting}

\subsubsection{Notification}\label{DpNAhdkSnHg7DcMraAzeH-2}
\\
The \texttt{Notification} class is another abstract class responsible for sending notifications, with the \texttt{SMSNotification} and \texttt{EmailNotification} classes as its children. The implementation of this class is given below:

\textbf{Notification and its derived classes}

\begin{lstlisting}[language=Java]
// Notification is an abstract class
public abstract class Notification {
    private int notificationId;
    // The Date data type represents and deals with both date and time.
    private Date createdOn;      
    private String content;

    public abstract void sendNotification(Person person);
}

class SMSNotification extends Notification {

    public void sendNotification(Person person) {
        // functionality 
    }
}

class EmailNotification extends Notification {

    public void sendNotification(Person person) {
        // functionality 
    }
}
\end{lstlisting}

\subsubsection{Room, room key, and room maintenance}\label{yWstzcn_6xVAdH6rz-InF}
\\
The \texttt{Room} class represents a room in the hotel. \texttt{RoomKey} is a class used to express the electronic key card, and \texttt{RoomMaintenance} is a class used to keep track of all the maintenance records for the rooms. The implementation of these classes is given below:

\textbf{The Room, RoomKey, and RoomMaintenance classes}

\begin{lstlisting}[language=Java]
public class Room {
    private String roomNumber;
    private RoomStyle style;
    private RoomStatus status;
    private double bookingPrice;
    private boolean isSmoking;
    private List<RoomKey> keys;
    private List<RoomHousekeeping> housekeepingLog;

    public boolean isRoomAvailable();
    public boolean checkin();
    public boolean checkout();
}

public class RoomKey {
    private String keyId;
    private String barcode;
    private Date issuedAt;
    private boolean isActive;
    private boolean isMaster;

    public boolean assignRoom(Room room);
}

public class RoomHousekeeping
{
    private String description;
    private Date startDatetime;
    private int duration;
    private Housekeeper housekeeper;

    public boolean addHousekeeping(Room room);
}
\end{lstlisting}

\subsubsection{Search and catalog}\label{_FM1DBRLxFVNgQAcR3p2f-2}
\\
\texttt{Search} is an interface, and the \texttt{Catalog} class implements the search interface to help in the room search. The code to perform this functionality is presented below:

\textbf{The Search interface and Catalog class}

\begin{lstlisting}[language=Java]
public interface Search {
    public static List<Room> search(RoomStyle style, Date date, int duration);
}

public class Catalog implements Search {
    private List<Room> rooms;

    public List<Room> search(RoomStyle style, Date date, int duration);
}
\end{lstlisting}

\subsubsection{Hotel and hotel branch}\label{tSGf5z-W4P90gYIxbJgl6}
\\
The \texttt{Hotel} class is the base class of the system that represents the hotel. The implementation of these classes is given below:

\textbf{The HotelBranch and Hotel classes}

\begin{lstlisting}[language=Java]
public class HotelBranch {
  private String name;
  private Address location;

  public List<Room> getRooms();
}

public class Hotel {
  private String name;
  private List<HotelBranch> locations;

  public boolean addLocation(HotelBranch location);
}
\end{lstlisting}

\subsection{Executable code: Hotel management system}\label{bqWeE4OCT4mvEnmDZsB6A}
\\
Below is a fully self-contained, runnable program in Java, C\#, C++, Python, and JavaScript that demonstrates the core workflows of the Hotel Management system. The main driver code of the system resides in the~\texttt{Driver.java},~\texttt{Driver.cs},~\texttt{Driver.py},~\texttt{Driver.cpp}, and~\texttt{Driver.js}~for all respective languages. You can click the ``Run'' button to execute the codes.

\subsubsection{What does this code show}\label{--bFT1y0UbtMW7RoYpx_K}
\\
This code simulates a basic hotel management system workflow involving hotel setup, room booking, notifications, and billing. It demonstrates how different components---such as hotels, branches, rooms, guests, bookings, and transactions---interact within the system in a real-world scenario.

\begin{itemize}
\item
\protect\phantomsection\label{CDd_uwl7OwdgD9u14hBpS}
\textbf{Scenario 1 (Hotel and room setup):} A hotel named ``Ocean View Resort'' has two branches: the \emph{Beachside} and \emph{City Center}. A catalog is initialized with two rooms---one standard and one deluxe. This forms the foundational setup for booking and operational activities.
\item
\protect\phantomsection\label{h5XEvKoDpyqrkxucaLXO5}
\textbf{Scenario 2 (Room search, booking, and notification):} The system searches for available ``STANDARD'' rooms. A guest named Alice Smith is created, and a booking is made for one of the available rooms. Email and SMS notifications are triggered and sent to the guest upon successful booking, showcasing customer communication capabilities.
\item
\protect\phantomsection\label{g7WsxyZRYs2JIKwSWX-JE}
\textbf{Scenario 3 (Payment and checkout process):} A cash transaction is initiated for the guest's stay, marking the billing step of the process. After the stay is complete, the room is checked out, and its status is updated. This scenario represents a guest's interaction with the hotel---from booking to payment and checkout.
\end{itemize}
\begin{quote}
\textbf{Hint: Try customizing the code!}\\
This code is designed for experimentation, not just reading. You can edit, extend, or modify any part of the example.
\end{quote}

\textbf{Hotel Management System}

\begin{lstlisting}[language=Java]
import java.util.Date;
import java.util.List;
import java.util.ArrayList;

public class Driver {
    public static void main(String[] args) {
        // ===============================
        // 🏨 Scenario 1: Hotel and Room Setup
        // ===============================
        System.out.println("=== Scenario 1: Hotel and Room Setup ===");

        // Step 1: Create the hotel
        Hotel hotel = new Hotel("Ocean View Resort");
        System.out.println("Hotel created: " + hotel.getName());

        // Step 2: Add branches to the hotel
        HotelBranch beachBranch = new HotelBranch("Beachside Branch", new Address());
        HotelBranch cityBranch = new HotelBranch("City Center Branch", new Address());

        hotel.addLocation(beachBranch);
        hotel.addLocation(cityBranch);
        System.out.println("Added " + hotel.getLocations().size() + " hotel branches.");

        // Step 3: Create and add rooms
        Room room101 = new Room("101", RoomStyle.STANDARD, RoomStatus.AVAILABLE, 120.0, false);
        Room room202 = new Room("202", RoomStyle.DELUXE, RoomStatus.AVAILABLE, 220.0, true);

        Catalog catalog = new Catalog();
        catalog.getRooms().add(room101);
        catalog.getRooms().add(room202);
        System.out.println("Room catalog initialized with " + catalog.getRooms().size() + " rooms.
");

        // ===============================
        // 🛏️ Scenario 2: Room Search, Booking & Notification
        // ===============================
        System.out.println("=== Scenario 2: Room Search, Booking & Notification ===");

        // Step 1: Search for STANDARD rooms
        System.out.println("Searching for available STANDARD rooms for 2 nights...");
        List<Room> availableRooms = catalog.search(RoomStyle.STANDARD, new Date(), 2);
        System.out.println("Available STANDARD rooms found: " + availableRooms.size());

        // Step 2: Guest creation and room booking
        Guest guest = new Guest();
        guest.setName("Alice Smith");
        guest.setEmail("alice.smith@example.com");
        guest.setPhone("987-654-3210");

        RoomBooking booking = new RoomBooking();
        booking.setReservationNumber("RS789123");
        booking.setRoom(room101);
        booking.setStartDate(new Date());
        booking.setDurationInDays(2);
        booking.setStatus(BookingStatus.CONFIRMED);

        guest.getBookings().add(booking);
        System.out.println("Booking confirmed for guest: " + guest.getName());
        System.out.println("Reservation Number: " + booking.getReservationNumber());

        // Step 3: Notifications
        EmailNotification email = new EmailNotification();
        email.setNotificationId(101);
        email.setContent("Hello " + guest.getName() + ", your room booking is confirmed!");
        email.setCreatedOn(new Date());
        email.sendNotification(guest);

        SMSNotification sms = new SMSNotification();
        sms.setNotificationId(102);
        sms.setContent("Hi " + guest.getName() + ", booking confirmed. Enjoy your stay!");
        sms.setCreatedOn(new Date());
        sms.sendNotification(guest);

        System.out.println();

        // ===============================
        // 💳 Scenario 3: Payment & Checkout Process
        // ===============================
        System.out.println("=== Scenario 3: Payment & Checkout Process ===");

        // Step 1: Initiate payment
        BillTransaction transaction = new CashTransaction();
        transaction.setAmount(room101.getBookingPrice());
        transaction.setCreationDate(new Date());
        transaction.setStatus(PaymentStatus.PENDING);
        transaction.initiateTransaction();

        System.out.println("Payment of $" + transaction.getAmount() + " initiated using cash.");
        System.out.println("Transaction status: " + transaction.getStatus());

        // Step 2: Room checkout
        room101.checkout();
        System.out.println("Room " + room101.getRoomNumber() + " has been successfully checked out.");
        System.out.println("Room status is now: " + room101.getStatus());

        System.out.println("
✅ All hotel operations completed successfully!");
    }
}
\end{lstlisting}

\subsection{Wrapping up}\label{_FM1DBRLxFVNgQAcR3p2f-3}
\\
In this chapter, we've explored the complete design of a hotel management system. We 've also shown that the system can be visualized using various UML diagrams and designed using object-oriented principles and design patterns.

% End of section content